\documentclass[11pt]{article}

\usepackage[margin=1in]{geometry}
\usepackage{amsmath,amssymb}
\usepackage{graphicx}
\usepackage{booktabs}
\usepackage[hidelinks]{hyperref}

\setlength{\parindent}{0pt}
\setlength{\parskip}{0.6\baselineskip}

\title{
U-Net with ResNet34 Encoder for Solar Filament Segmentation:
Training at 1024\,px with Component-wise Instance Filtering
}

\author{
Mauro Frank Lima de Lima\\
\small Independent Researcher, Manaus, Amazonas, Brazil\\
\small \texttt{mauro.frank.lima@outlook.com}
}

\date{}

\begin{document}

\maketitle

\begin{abstract}
Solar filaments are elongated structures observed in H-alpha solar
images. Their segmentation is challenging because of their thin
morphology, irregular boundaries, low contrast, and the presence of
other structures and noise in the solar atmosphere.

This work presents a pixel-level segmentation pipeline based on
U-Net with a ResNet34 encoder, trained using MAGFiLO H-alpha images.
The inference pipeline combines eight-view test-time augmentation
(TTA), probability-map averaging, and component-wise filtering to
generate individual filament predictions.

The input images are normalized using the 1st and 99th percentiles,
converted to grayscale, and resized from 2048$\times$2048 to
512$\times$512 pixels. During inference, predictions are restored
to the original image resolution and processed using probability
thresholds and morphological closing.

The submitted inference pipeline (v13) generated 1790 predicted
instances across 178 test images. The model achieved a validation
Dice of 0.627 and a public leaderboard score of 0.28.
\end{abstract}


\section{Introduction}

Solar filaments are relatively cool and dense plasma structures
suspended in the solar corona by magnetic fields. They are commonly
observed as dark elongated features against the brighter solar
chromosphere in H-alpha images.

Their morphology and evolution are relevant to the study of solar
activity and space weather. Reliable identification of these
structures can support the analysis of their spatial distribution,
morphological properties, and evolution.

Automatic filament segmentation is challenging because filaments
may exhibit thin structures, irregular boundaries, branches,
and complex shapes. Their appearance can also be affected by
image noise, variations in contrast, sunspots, and other
chromospheric structures.

These characteristics make pixel-level segmentation more
appropriate than image-level classification. In addition,
distinguishing individual filament instances requires separating
connected structures rather than producing only a single
foreground mask.

This work investigates a segmentation pipeline based on U-Net
with a ResNet34 encoder. Test-time augmentation and component-wise
probability filtering are incorporated into the inference stage
to generate individual filament predictions.

Panoptic Quality (PQ) is a metric commonly used for evaluating
instance-level segmentation. It combines segmentation quality
and recognition quality. However, PQ values are not reported
in this work because the corresponding validation evaluation
was not completed in the recorded experimental run.


\section{Data and Preprocessing}

The experiments use H-alpha solar images from the MAGFiLO dataset
\cite{magfilo}.

The original images have a resolution of 2048$\times$2048 pixels.
The preprocessing pipeline consists of the following operations:

\begin{enumerate}
    \item Convert the input image to grayscale.
    \item Normalize image intensities using the 1st and 99th
          percentiles.
    \item Resize the image from 2048$\times$2048 to
          512$\times$512 pixels.
    \item Rasterize the expert annotations from COCO polygon
          representations into binary segmentation masks.
\end{enumerate}

The dataset is divided into training and validation subsets
using an 80/20 split by image file. This strategy is intended
to prevent images associated with the same file from being
distributed across both subsets.

Annotations belonging to category~4 (Ambiguous) are excluded
from training. The remaining three categories (Left-bearing,
Right-bearing, and Unidentifiable) are treated as a single
foreground class.


\section{Model and Training}

The segmentation architecture is a U-Net with a ResNet34
encoder, implemented using the
\texttt{segmentation\_models\_pytorch} library.

The network receives single-channel grayscale images and
produces pixel-level segmentation predictions.

The training configuration is as follows:

\begin{itemize}
    \item Architecture: U-Net with ResNet34 encoder.
    \item Encoder weights initialised from ImageNet pre-training.
    \item Input channels: 1 (grayscale).
    \item Input resolution: 1024$\times$1024 pixels (random
          512$\times$512 crops during training).
    \item Loss function: Dice loss (smooth\,=\,1) combined with
          binary cross-entropy, positive class weight 20.
    \item Optimizer: AdamW, learning rate $5\times10^{-4}$.
    \item Batch size: 4.
    \item Training duration: 35 epochs (best at epoch~33).
    \item Mixed-precision training (fp16 autocast + GradScaler).
    \item Gradient clipping (max norm 1.0).
\end{itemize}

On-the-fly data augmentation includes random horizontal flip,
vertical flip, and 90-degree rotation (uniform over
$\{0^\circ, 90^\circ, 180^\circ, 270^\circ\}$).

The model checkpoint used by the inference pipeline is
\texttt{weights/best.pt} (float32, 98\,MB), available in the
public repository (v13 weights).


\section{Inference and Post-processing}

\paragraph{Test-time augmentation.}
The inference pipeline applies eight-view TTA, consisting of
four rotations ($0^\circ$, $90^\circ$, $180^\circ$, $270^\circ$)
combined with horizontal flip on/off. The model predicts a
probability map for each transformed input. The corresponding
inverse transformations are applied so that all predictions
return to the original spatial orientation. The eight probability
maps are averaged pixel-wise to obtain the final probability map.

\paragraph{Upsampling.}
The averaged 512$\times$512 probability map is upsampled to
2048$\times$2048 with bilinear interpolation.

\paragraph{Component-wise filtering.}
A seed mask is obtained by thresholding at $t_{\mathrm{lo}}$,
followed by morphological closing with an elliptical kernel
(radius 5\,px). Connected components are labelled; a component
is retained as a filament instance only if all three conditions
hold:
\begin{itemize}
    \item $\max(\mathrm{prob}) > t_{\mathrm{hi}}$ (peak confidence),
    \item $\mathrm{mean}(\mathrm{prob}) \geq \mathrm{min\_meanp}$
          (average confidence),
    \item $\mathrm{area} \geq \mathrm{min\_area}$ (minimum size).
\end{itemize}
This strategy is inspired by~\cite{bahrami2026}, without the
gradient-boosting post-filter described there.

\paragraph{Threshold configuration.}
The thresholds used in the submitted inference are:
$t_{\mathrm{lo}} = 0.35$,
$t_{\mathrm{hi}} = 0.50$,
$\mathrm{min\_meanp} = 0.30$,
$\mathrm{min\_area} = 50$\,pixels.

Each retained component is encoded separately using
run-length encoding (RLE), producing one instance per
predicted filament.


\section{Results}

The v13 inference pipeline generated 1790 predicted instances
across 178 test images. The predictions were encoded using COCO
binary RLE and stored in \texttt{submission\_v13.csv}.

\begin{table}[ht]
\centering
\caption{Summary of the submitted inference output (v13).}
\begin{tabular}{lc}
\toprule
Quantity & Result \\
\midrule
Test images processed       & 178 \\
Predicted instances         & 1790 \\
TTA views                   & 8   \\
Training input resolution   & $1024\times1024$ (crop $512\times512$) \\
Output resolution           & $2048\times2048$ \\
Validation Dice (semantic)  & 0.627 \\
Public leaderboard score    & 0.28 \\
Training epochs             & 35 (best at epoch 33) \\
\bottomrule
\end{tabular}
\end{table}

\begin{table}[ht]
\centering
\caption{Comparison of submitted versions.}
\begin{tabular}{lccc}
\toprule
Version & Epochs & Val Dice & Public LB \\
\midrule
v12 (15\,ep, 512\,px) & 15 & --- & 0.00 \\
v13 (35\,ep, 1024\,px) & 35 & 0.627 & 0.28 \\
\bottomrule
\end{tabular}
\end{table}

The v12 submission scored 0.00 on the public leaderboard. Post-hoc
diagnosis identified that the default thresholds in the inference
script ($t_{\mathrm{hi}}=0.99$, $\mathrm{min\_meanp}=0.85$) were
too restrictive, suppressing all predictions. The v13 retraining
also increased the input resolution from 512 to 1024\,px and
extended training to 35 epochs, yielding a validation Dice of 0.627
and a public score of 0.28.


\section{Reproducibility}

The source code, trained weights (\texttt{weights/best.pt},
float32), and inference notebook (\texttt{pipeline.ipynb}) are
available at:

\begin{center}
\url{https://github.com/maurolima-sys/solar-filament-segmentation}
\end{center}

To reproduce: install dependencies with
\texttt{pip install -r requirements.txt} and run
\texttt{pipeline.ipynb}, which loads \texttt{weights/best.pt}
without retraining. GPU used for training: NVIDIA T4 (Kaggle).


\begin{thebibliography}{9}

\bibitem{magfilo}
A.~Ahmadzadeh, D.~J.~Kempton, Q.~Li, A.~A.~Pevtsov,
``MAGFiLO: Manually Annotated GONG Filaments from H-Alpha
Observations,''
\textit{Scientific Data}, 2024.
\url{https://doi.org/10.1038/s41597-024-03876-y}

\bibitem{pq}
A.~Kirillov et al.,
``Panoptic Segmentation,''
\textit{Proc.\ IEEE/CVF CVPR}, 2019.
\url{https://doi.org/10.1109/CVPR.2019.00963}

\bibitem{bahrami2026}
D.~Bahrami,
``Solar Filament U-Net: OOF 0.45 vs Public 0.39,''
Kaggle notebook, 2026.
\url{https://www.kaggle.com/code/dariushafshar/solar-filament-u-net-oof-0-45-vs-public-0-39}

\end{thebibliography}

\end{document}
